\documentclass[12pt,a4paper]{article}

\usepackage[utf8]{inputenc}
\usepackage[T1]{fontenc}
\usepackage[norsk]{babel}
\usepackage{geometry}
\usepackage{microtype}
\usepackage{setspace}
\usepackage{booktabs}
\usepackage{longtable}
\usepackage{array}
\usepackage{enumitem}
\usepackage{xcolor}
\usepackage{hyperref}

\geometry{margin=2.6cm}
\onehalfspacing
\setlist[itemize]{leftmargin=*, itemsep=0.2em}
\setlist[enumerate]{leftmargin=*, itemsep=0.2em}

\hypersetup{
  colorlinks=true,
  linkcolor=black,
  urlcolor=blue,
  citecolor=black
}

\newcommand{\todo}[1]{\par\noindent\textcolor{red!70!black}{\textbf{TODO:} #1}\par}
\newcommand{\claimlabel}[1]{\textit{#1}}

\title{Coop App, Lobyco og digital strategi\\
\large Arbeidsutkast i LaTeX}
\author{IT-strategioppgave}
\date{4. mai 2026}

\begin{document}

\maketitle

\begin{center}
\begin{tabular}{ll}
\toprule
Status & Arbeidsutkast, ikke ferdig rapport \\
Språk & Norsk \\
Rådgivningsaktør & Martin Hasgard Olesen, Direktør for Kommunikation, Marketing og Digital \\
Perspektiver & DVC Framework og Digital Business Model-perspektivet \\
Strategisk stance & Hybrid governance \\
\bottomrule
\end{tabular}
\end{center}

\tableofcontents
\newpage

\section*{Arbeidsstatus og agentmodell}
\addcontentsline{toc}{section}{Arbeidsstatus og agentmodell}

De sentrale prosjektvalgene er låst: oppgaven rådgir Martin Hasgard Olesen, bruker DVC Framework og Digital Business Model-perspektivet, skrives på norsk, og bygger mot en hybrid governance-stance for Coop App og Lobyco. De valgte anbefalingene handler om å redefinere app-suksess rundt forretningsmodell-fit, utvikle kjedespesifikke app-value propositions innenfor én styrt plattform, gjennomgå gamification og engagement-funksjoner, styrke medlemsrelasjonen uten å forveksle lojalitet med lønnsomhet, og teste kjedenivå-initiativer gjennom piloter før skalering.

Agentmodellen brukes eksplisitt i dette utkastet:

\begin{itemize}
  \item \textbf{Exam Architect:} holder strukturen tett på oppgavekravet om analyse først og anbefalinger etterpå.
  \item \textbf{Evidence Extractor:} kobler sentrale påstander til kildepakke og evidensfiler.
  \item \textbf{Fact Verification Agent:} kontrollerer at konkrete fakta, tall, datoer, navn og regnskapsbegreper stemmer med kilde- og verifikasjonsfilene før de brukes som fakta.
  \item \textbf{Source Critic:} markerer hvor Lobyco-, Playable- og Shortcut-påstander må behandles forsiktig.
  \item \textbf{Competitor \& Market Agent:} bruker konkurrentevidens til å vise digital/business-model fit, ikke kausal app-effekt.
  \item \textbf{Theory Application Agent:} bruker bare DVC og Digital Business Model som fulle teoriperspektiver.
  \item \textbf{Recommendation Agent:} sikrer at senere anbefalinger følger fra analysen.
  \item \textbf{Red Team Agent:} passer på at Lobyco, gamification, retail media og lojalitet ikke blir behandlet som dokumentert lønnsomhet.
\end{itemize}

Dette LaTeX-utkastet bygger videre på \texttt{draft\_v2.md}. Det inneholder nå sammenhengende rapporttekst for innledning, metode, casebakgrunn, kildekritikk, teorivalg/anvendelse, strategisk situasjonsanalyse, strategisk stance, strategiske anbefalinger, begrensninger, konklusjon og AI-bruk. Referanseseksjonen er fortsatt en verifiseringsliste, fordi prosjektet ikke skal finne på referanser, sidetall eller bibliografiske detaljer før kildene er kontrollert.

\section*{Fact Verification Agent: kontroll før skriving}
\addcontentsline{toc}{section}{Fact Verification Agent: kontroll før skriving}

Dette utkastet bruker et eksplisitt faktasjekktrinn før nye påstander skrives som rapporttekst. Hensikten er å skille mellom verifiserte fakta, kildeclaims, marketingclaims, antagelser og analytiske inferenser.

\subsection*{Faktapunkter brukt som relativt sterke fakta}

\begin{longtable}{p{0.34\textwidth} p{0.29\textwidth} p{0.29\textwidth}}
\toprule
\textbf{Faktapunkt} & \textbf{Status i prosjektgrunnlaget} & \textbf{Bruk i utkastet} \\
\midrule
\endfirsthead
\toprule
\textbf{Faktapunkt} & \textbf{Status i prosjektgrunnlaget} & \textbf{Bruk i utkastet} \\
\midrule
\endhead
Coop Danmark A/S rapporterer 32,565 mrd. DKK i nettoomsetning i 2025 & Verifisert i årsrapport-/kildepakke & Brukes som fysisk dagligvare- og skalaindikator \\
Coop rapporterer 44,756 mrd. DKK i nettoomsetning inkludert brugsforeninger & Verifisert i årsrapport-/kildepakke & Brukes som bredere Coop-family retail scale \\
Coop Danmark A/S hadde 536 butikker, 900 inkludert brugsforeninger & Verifisert i årsrapport-/kildepakke & Brukes som grunnlag for bricks-and-mortar-kjerne \\
EBITDA før eiendomsgevinster var 313 mill. DKK, mens EBIT var -215 mill. DKK & Verifisert, men krever presise regnskapsbegreper & Brukes som investeringspress, ikke app-årsak \\
Coop.dk MAD ble lukket i 2023 og Coop.dk-webshop 31. januar 2025 & Verifisert i kildepakke, årsak må formuleres forsiktig & Brukes som kanal- og forretningsmodellkontekst \\
Coop har mer enn 2 millioner medlemmer/medeiere & Offisiell Coop-claim/fakta i kildepakke & Brukes som medlems- og relasjonskontekst \\
\bottomrule
\end{longtable}

\subsection*{Påstander behandlet som claims eller marketingclaims}

\begin{longtable}{p{0.34\textwidth} p{0.29\textwidth} p{0.29\textwidth}}
\toprule
\textbf{Påstand} & \textbf{Status i prosjektgrunnlaget} & \textbf{Bruk i utkastet} \\
\midrule
\endfirsthead
\toprule
\textbf{Påstand} & \textbf{Status i prosjektgrunnlaget} & \textbf{Bruk i utkastet} \\
\midrule
\endhead
Lobyco oppgir 1,8m+ appbrukere, ca. 25 \% household reach og ca. 880 000 MAU & Leverandør-/marketingclaim & Brukes som reach/adoption-indikator, ikke lønnsomhetsbevis \\
Lobyco hevder at appbrukere handler oftere & Marketingclaim med kausalitets- og seleksjonsrisiko & Brukes kun som forsiktig claim \\
Playable rapporterer spill-, engagement- og innløsningsdata & Leverandør-/marketingclaim & Brukes som mulig mekanisme, ikke nettoeffekt \\
Shortcut beskriver apputviklingshistorikk, downloads og OEM-/plattformlogikk & Utvikler-/leverandørclaim & Brukes som historikk og kapabilitetsclaim \\
Konkurrentevidensen antyder digital/business-model fit & Analytisk inferens fra kildepakke & Brukes som kontekst, ikke kausal proof \\
\bottomrule
\end{longtable}

\subsection*{TODO-punkter for senere faktasjekk}

\begin{itemize}
  \item Sett inn endelige APA-referanser til årsrapport, Coop-sider, Lobyco, Playable, Shortcut og teori.
  \item Sjekk endelig kildeordlyd for årsakene til Coop.dk MAD-/Coop.dk-webshop-lukking.
  \item Sjekk definisjoner, datoer og målegrunnlag for Lobyco-tallene ``brukere'', ``MAU'' og ``household reach''.
  \item Avklar 232 mill. DKK loss-claim fra case-/prosjektkontekst opp mot -215 mill. DKK EBIT før finalrapport.
\end{itemize}

\section{Innledning}

Coop står i 2026 overfor en strategisk diskusjon om hvilken rolle Coop App og Lobyco bør ha i konsernets videre utvikling. Spørsmålet er ikke bare om Coop skal investere mer eller mindre i appen. Det handler også om hvordan en kunderettet digital kanal skal forstås i en virksomhet der den fysiske butikkstrukturen fortsatt utgjør kjernen i forretningsmodellen.

Casebeskrivelsen peker på flere mulige veivalg: Coop kan redusere investeringene i appen, eventuelt afhende Lobyco, styrke den videre apputviklingen, utvikle mer forskjellige app-løsninger for de ulike kjedene, eller velge en annen mellomposisjon. Disse alternativene må vurderes mot Coops økonomiske situasjon, kjedeportefølje, medlemsmodell og fysiske butikkgrunnlag.

Denne oppgaven rådgir Martin Hasgard Olesen, Direktør for Kommunikation, Marketing og Digital. Valget av Martin gjør det naturlig å analysere Coop App som en kunderettet digital kanal, men også som et styrings- og prioriteringsspørsmål. Martin må kunne delta i direksjonsdiskusjonen med et grunnlag som skiller mellom digital aktivitet, kundeverdi og faktisk forretningsmessig verdi.

Oppgavens overordnede problemstilling er:

\begin{quote}
Hvordan bør Martin Hasgard Olesen forholde seg til fremtiden for Coop App og Lobyco i den forventede direksjonsdiskusjonen, gitt Coops økonomiske press, fysiske butikkjerne, kjedeforskjeller og usikre evidensgrunnlag for appens økonomiske verdi?
\end{quote}

Argumentet i oppgaven er at Martin ikke bør anbefale en enkel styrk-, kutt- eller selg-logikk. I stedet bør han argumentere for hybrid governance: Coop bør beholde Coop App og Lobyco som mulige strategiske digitale kapabiliteter, men videre investeringer bør knyttes til butikkøkonomi, kjedespesifikk relevans, kundetillit og tydelige KPI-er. Med andre ord bør appen ikke måles primært på reach, kampanjeaktivitet eller engagement, men på om den styrker Coops butikkbaserte forretningsmodell.

Oppgaven er strukturert slik: Først presenteres metode, avgrensninger og sentrale antagelser. Deretter gis en kort casebakgrunn om Coop, Coop App og Lobyco. Videre diskuteres kildegrunnlaget kritisk, før DVC Framework og Digital Business Model-perspektivet anvendes i en strategisk situasjonsanalyse. På bakgrunn av analysen utvikles fem strategiske anbefalinger til Martin Hasgard Olesen.

\section{Metode, avgrensning og antagelser}

Oppgaven anvender to selvvalgte perspektiver fra undervisningen: DVC Framework og Digital Business Model-perspektivet. Perspektivene brukes ikke som brede, generelle modeller for digital transformasjon. De brukes selektivt for å analysere den delen av Coop-casen som handler om Coop App, Lobyco, medlemsrelasjoner, kjedespesifikk digital relevans og forretningsmodellmessig verdi.

DVC Framework brukes til å analysere hvilke former for digital verdi Coop App kan skape. Det gjelder særlig digitale kundeopplevelser, relasjoner, relevans, digital infrastruktur og digitale output-mål. Perspektivet er nyttig fordi Coop App kan skape verdi på flere måter: gjennom tilbud, medlemsfordeler, personalisering, butikkstøtte, gamification og kontakt med medlemmene. Samtidig gjør DVC det mulig å skille mellom synlig digital aktivitet og dypere kundeverdi.

Digital Business Model-perspektivet brukes til å vurdere om denne digitale verdien passer med Coops forretningsmodell. Det sentrale spørsmålet er om appen og Lobyco bidrar til value proposition, value creation og value capture i en virksomhet som fortsatt er sterkt forankret i fysisk butikkhandel. Perspektivet brukes også til å analysere Lobyco som mulig kapabilitet, appen som kanal og behovet for governance og KPI-er.

Analysen bygger på et sammensatt kildegrunnlag. Offisielle Coop-kilder og årsrapportmateriale brukes som sterkest grunnlag for finansielle og strukturelle påstander. Coop rapporterer blant annet 32,565 mrd. DKK i nettoomsetning for 2025, EBITDA før eiendomsgevinster på 313 mill. DKK og et EBIT-underskudd på 215 mill. DKK. Samtidig viser kildegrunnlaget at Coop er en stor fysisk butikkaktør med 536 Coop Danmark-butikker og 900 butikker inkludert brugsforeninger. Disse tallene brukes som bakgrunn for å forstå investeringspress og butikkbasert strategisk kjerne. Tallene må presenteres med presise regnskapsbegreper og ikke blandes med andre loss- eller resultatmål.

Kilder om Coop App, Lobyco, Playable og Shortcut brukes mer forsiktig. Lobyco oppgir for eksempel at Coop App har mer enn 1,8 millioner brukere og omtrent 25 prosent household reach, men dette er reach- og adoption-data, ikke dokumentasjon på lønnsomhet. På samme måte kan Playable-kilder indikere engagement og mulig leverandørfinansiert kampanjeverdi, men de kan ikke alene bevise netto økonomisk effekt for Coop. Denne kildekritikken er viktig fordi oppgaven skal vise en selvstendig og autoritativ vurdering, ikke bare gjengi leverandørclaims.

Konkurrentmateriale brukes som kontekst, ikke som direkte bevis. Salling, REMA 1000, Lidl, Dagrofa/MENY og Nemlig viser ulike måter digitale løsninger kan knyttes til kundeløfte, butikkdrift og forretningsmodell. Denne evidensen brukes til å formulere en forsiktig inferens: digitale verktøy ser ut til å være mest relevante når de passer med virksomhetens kjerneposisjon og driftslogikk. Den brukes ikke til å hevde at konkurrenter lykkes fordi de har apper.

Oppgaven bygger på fire sentrale antagelser:

\begin{enumerate}
  \item Coop App behandles som en butikkstøttende digital kanal, ikke som en erstatning for en full e-commerce-strategi.
  \item Appbruk, kampanjeaktivitet og engagement behandles som indikatorer, ikke som bevis på profitabilitet.
  \item Lobyco behandles som et åpent governance- og kapabilitetsspørsmål, ikke som noe som automatisk bør beholdes eller selges.
  \item Kjedespesifikk digital utvikling bør forstås som modulering innenfor en styrt plattform, ikke nødvendigvis som utvikling av helt separate apper.
\end{enumerate}

Disse avgrensningene betyr at oppgaven ikke forsøker å bevise appens netto økonomiske bidrag. Det offentlige kildegrunnlaget gir ikke tilstrekkelig data om margin, kontrollgrupper, incrementalitet, appkostnader, retail media-inntekter eller Lobycos økonomiske bidrag. I stedet analyserer oppgaven hvilke styringsspørsmål Martin bør reise, og hvilke anbefalinger som følger dersom Coop ønsker at appen skal støtte den butikkbaserte forretningsmodellen mer presist.

\section{Casebakgrunn: Coop, Coop App og Lobyco}

Coop Danmark bør i denne oppgaven først forstås som en fysisk butikkbasert dagligvareaktør. Årsrapport- og kildegrunnlaget viser at Coop Danmark A/S rapporterer 32,565 mrd. DKK i nettoomsetning i 2025, og 44,756 mrd. DKK når brugsforeningene inkluderes. Samme kildegrunnlag viser 536 Coop Danmark-butikker og 900 butikker inkludert brugsforeninger. Disse tallene gjør det lite presist å analysere Coop som en digital plattformvirksomhet i snever forstand. Coop har digitale kundekanaler og digitale kapabiliteter, men den økonomiske hovedarenaen er fortsatt dagligvarehandel gjennom fysiske butikker. \todo{Sett inn endelig APA-referanse til Coop Danmark A/S årsrapport 2025.}

Dette er viktig fordi caseproblemet handler om fremtiden for Coop App og Lobyco. Dersom Coop hadde vært en primært online dagligvareaktør, ville appen naturlig kunne analyseres som hovedkanalen for salg og kundeopplevelse. Kildegrunnlaget peker i motsatt retning. Coop.dk MAD ble lukket i 2023, og Coop.dk-webshop ble lukket 31. januar 2025. Kildene indikerer at økonomi og fysisk butikkfokus var sentrale forklaringer, men årsaksformuleringen må holdes forsiktig til den endelige case- og kildeordlyden er sjekket. Den strategiske konsekvensen er likevel tydelig nok for analysen: Coop App bør ikke behandles som en ny e-commerce-satsing som skal erstatte butikkene, men som et digitalt lag rundt butikkbesøk, medlemsrelasjon, tilbud, betaling, kvitteringer, personalisering og lokal relevans.

Coops kjedeportefølje gjør dette mer krevende. Coop består ikke av ett enkelt butikkonsept med ett kundeløfte. Kildegrunnlaget peker på flere kjeder og posisjoner, blant annet Kvickly, SuperBrugsen, Brugsen og 365discount. Disse kjedene har ulike kundesituasjoner, prisforventninger og handleoppdrag. En kunde som bruker 365discount, kan primært lete etter lav pris og enkelhet, mens kunder i andre Coop-kjeder kan legge mer vekt på bredde, ferskvarer, lokal tilhørighet, inspirasjon eller medlemsfordeler. Dette betyr ikke at hver kjede nødvendigvis bør ha en egen app. Det betyr derimot at én felles app-logikk kan være strategisk svak dersom den ikke oppleves relevant på tvers av kjedene.

Coop App fremstår i de offentlige kildene som en medlems-, lojalitets- og butikkstøttende kanal. Coop beskriver appfunksjoner som personlige tilbud, bonus, Scan \& Betal, digitale kvitteringer, handleliste, spill, klimaavtrykk og lokal butikkkommunikasjon. Dette viser at appen har flere mulige verdilogikker samtidig. Den kan redusere friksjon i butikkreisen, gjøre tilbud mer relevante, støtte medlemsrelasjonen og skape nye kontaktpunkter mellom Coop, medlemmet, butikken og leverandører. Samtidig er denne bredden også en strategisk risiko: hvis appen blir målt på for mange aktivitetsmål uten tydelig kobling til butikkøkonomi, kan digital aktivitet bli forvekslet med forretningsmessig verdi.

Medlemsmodellen er et særtrekk ved Coop-casen. Coop oppgir mer enn 2 millioner medlemmer/medeiere, og dette gir appen en annen rolle enn en ordinær rabatt- eller kampanjeapp. Appen kan brukes til å forsterke medlemsidentitet, personlig relevans og kommunikasjon mellom Coop og kundene. Likevel må medlemsverdi og profitabilitet holdes analytisk atskilt. En sterk medlemsrelasjon kan være strategisk verdifull selv om den ikke umiddelbart kan måles i kortsiktig margin. Men det er heller ikke riktig å anta at appbrukere er mer lønnsomme bare fordi de er mer aktive, mer lojale eller mer synlige i digitale data. Appbrukere kan allerede være blant Coops mest lojale kunder, og dermed kan observerte forskjeller mellom appbrukere og andre kunder være preget av seleksjonsbias.

Lobyco er en sentral del av casebakgrunnen fordi selskapet gjør Coop App-spørsmålet til mer enn et vanlig kanalspørsmål. Kildene beskriver Lobyco som knyttet til Coop App, multi-banner loyalty og en mulig plattform- eller OEM-logikk. Dette kan indikere at Coop har bygget en digital kapabilitet som kan ha verdi utover enkeltfunksjoner i appen. For Martin Hasgard Olesen er dette strategisk relevant: Lobyco kan forstås som digital infrastruktur, som kompetansemiljø, som mulig retail media- og lojalitetsplattform, eller som et investeringsobjekt som må vurderes opp mot andre konsernprioriteringer.

Samtidig er Lobyco-kildene ikke nøytrale. Lobyco oppgir at Coop App har mer enn 1,8 millioner brukere, omtrent 25 prosent household reach og rundt 880 000 månedlige aktive brukere. Dette viser en betydelig rapportert digital rekkevidde, men det er fortsatt leverandør-/marketingrapportert evidens. Det dokumenterer ikke i seg selv netto økonomisk bidrag, at Lobyco gir Coop et varig konkurransefortrinn, eller at appen skaper inkrementell handlefrekvens. Lobyco hevder også at appbrukere handler oftere, men dette bør behandles som en korrelasjons- eller marketingclaim, ikke som dokumentasjon på kausal effekt. \todo{Sjekk endelig ordlyd, dato og definisjon av ``bruker'', ``MAU'' og ``household reach'' i Lobyco-kildene.}

Playable- og Shortcut-kildene supplerer dette bildet, men de må brukes med samme forsiktighet. Playable beskriver gamification, premier, spilldeltakelse, innløsning i butikk og mulig retail media-/supplier-funded campaign-logikk. Dette kan brukes til å vise hvordan Coop App kan skape engagement og leverandørfinansiert aktivitet. Det kan ikke brukes som bevis for netto lønnsomhet, særlig fordi kampanjekostnader, margin, kontrollgrupper og langsiktig kundeatferd ikke er dokumentert i det offentlige kildegrunnlaget. Shortcut beskriver apputviklingshistorikk, downloads, daglige brukere og OEM-/plattformlogikk, men også dette er en utvikler-/leverandørfortelling snarere enn en uavhengig strategisk evaluering.

Coops økonomiske situasjon gjør kildebruken særlig viktig. Coop rapporterer positiv EBITDA før eiendomsgevinster på 313 mill. DKK, men også EBIT på -215 mill. DKK. Prosjektet inneholder dessuten et separat 232 mill. DKK loss-claim fra case-/prosjektkonteksten. Disse må ikke blandes sammen. I rapporten bør de finansielle tallene brukes til ett avgrenset poeng: Coop opererer under investeringspress, og derfor må app- og Lobyco-investeringer vurderes mot tydelige forretningsmodellkriterier. Tallene bør ikke brukes til å antyde at appen forklarer underskuddet, eller at kutt i appen automatisk vil løse Coops økonomiske problemer.

Konkurrentkonteksten forsterker dette poenget. Salling, REMA, Lidl, Dagrofa/MENY og Nemlig viser ikke at dagligvareaktører vinner fordi de har apper. Den sterkere og tryggere tolkningen er at digitale verktøy ser ut til å skape mer verdi når de passer med aktørens kundeløfte og driftsmodell. Salling kan koble digitale løsninger til skala og butikkflyt, Lidl Plus til kuponger og discount-logikk, REMA til enkel butikkutility, Dagrofa/MENY til mat- og lokal identitet, og Nemlig til en mer rendyrket online-grocery-modell med egne økonomiske utfordringer. For Coop betyr dette at spørsmålet ikke er om appen er ``digital nok'', men om den er relevant nok for Coops ulike kjeder og fysisk forankrede forretningsmodell.

Samlet peker casebakgrunnen mot én grunnleggende spenning: Coop har en stor fysisk butikk- og medlemsbase, en synlig app og en mulig digital kapabilitet i Lobyco, men det offentlige kildegrunnlaget dokumenterer ikke om disse digitale ressursene gir tilstrekkelig butikkøkonomisk og strategisk effekt. Derfor bør den videre analysen ikke starte med en normativ antagelse om at appen bør styrkes, kuttes eller selges. Den bør starte med spørsmålet om hvordan appen og Lobyco kan styres slik at digital kundeverdi faktisk passer med Coops forretningsmodell.

\section{Kildekritikk}

Kildegrunnlaget kan støtte en analyse av mulig appverdi, men det kan ikke bevise appens lønnsomhet eller kausale effekt på kundeatferd. Dette er et avgjørende premiss for oppgaven, fordi caseproblemet lett kan føre til to forenklede argumenter: enten at Coop App bør styrkes fordi den har mange brukere, eller at den bør reduseres fordi Coop er økonomisk presset. Begge slutninger er for bastante dersom de ikke kobles til tydelig evidens om kundeverdi, butikkøkonomi, kostnader og value capture.

De sterkeste kildene i materialet er Coop årsrapport og offisielle Coop-kilder. De kan brukes til å dokumentere rapporterte regnskapsstørrelser, butikkantall, butikkstruktur, medlemsmodell og offisielle beskrivelser av appfunksjoner. Også disse kildene er selskapets egen fremstilling, men for regnskapsmessige og strukturelle fakta er de klart sterkere enn leverandør- og kampanjekilder. Ved bruk av finansielle tall må rapporten bevare eksakte begreper: nettoomsetning, EBITDA før eiendomsgevinster, EBIT og eventuelle loss- eller årsresultatmål må ikke blandes.

Lobyco-kildene har en annen rolle. De er nyttige fordi de viser hvordan Lobyco selv fremstiller Coop App, multi-banner loyalty, appreach, månedlig aktivitet og mulig plattformlogikk. Men Lobyco er en kommersielt interessert kilde og bør derfor ikke brukes som nøytralt bevis for økonomisk effekt. Når Lobyco oppgir mer enn 1,8 millioner appbrukere og omtrent 25 prosent household reach, kan rapporten bruke dette som indikasjon på adoption og reach. Den bør samtidig eksplisitt si at dette ikke dokumenterer lønnsomhet, inkrementell handlefrekvens eller varig konkurransefortrinn.

Det samme gjelder Playable. Playable-kildene kan vise at gamification og supplier-funded campaigns er en del av mulig app-logikk. De kan også indikere at spill, premier og innløsning i butikk kan skape engagement og mulig kampanjeverdi. Kilden er likevel en leverandørcase. Den har en interesse i å vise at gamification fungerer. Derfor bør tall om spilldeltakelse, kurvverdi eller innløste premier presenteres som rapporterte kampanje- og engagement-tall, ikke som dokumentasjon på Coops netto økonomiske verdi. Et trygt resonnement er at gamification kan være en testbar mekanisme for kundeopplevelse og retail media, men at den må vurderes mot kostnader, margin, kundetillit og faktisk butikkadferd.

Shortcut-kilden bør hovedsakelig brukes til historikk og kapabilitetsbeskrivelse. Den kan støtte at Coop App har vært et betydelig apputviklingsprosjekt og at løsningen har hatt en utviklingshistorie knyttet til OEM- eller plattformlogikk. Den bør ikke brukes til å konkludere om appstrategien er vellykket i dag. Teknisk leveranse, downloads og daglige brukere er ikke det samme som strategisk verdi, og eldre utviklingstall kan være mindre relevante for Coops 2026-situasjon.

Konkurrentmaterialet er nyttig, men må avgrenses. Salling, REMA, Lidl, Dagrofa/MENY og Nemlig kan brukes til å vise mønstre i hvordan digitale løsninger kan passe med ulike forretningsmodeller. Salling-kilder kan dokumentere skala og digitale butikkfunksjoner, Lidl-kilder kan vise hvordan en app kan støtte kupong- og discount-logikk, og Dagrofa/MENY-kilder kan vise en annen lojalitets- og identitetslogikk. Dette gir relevant kontekst for Coop. Likevel kan rapporten ikke hevde at konkurrentenes økonomiske resultater skyldes appene. Konkurrentevidensen antyder digital fit, ikke direkte årsakssammenheng.

App-store- og scraping-relatert materiale bør brukes enda mer forsiktig. Prosjektets Google Play-scrape for Coop App feilet og skal ikke brukes til nøkkelpåstander om rating, downloads, appfunksjoner eller adoption. Lokale scraping-synteser kan være nyttige som arbeidsnotater, men finalrapportens sentrale claims bør så langt som mulig spores til originalkilder, årsrapporter, offisielle sider eller klart identifiserte leverandørkilder.

Den viktigste analytiske kildekritikken kan oppsummeres i fire nivåer:

\begin{quote}
Adoption/engagement $\rightarrow$ kundeatferd $\rightarrow$ butikkøkonomi $\rightarrow$ profitabilitet/strategisk fordel
\end{quote}

Mange av kildene dokumenterer eller påstår noe om de første nivåene. Lobyco og Playable kan si noe om reach, appbruk, kampanjedeltakelse og engagement. Coop kan beskrive appfunksjoner og medlemsfordeler. Det offentlige materialet sier derimot langt mindre om de senere nivåene: inkrementell atferd, margin, store traffic, basket effect, retail media-nettoinntekt, kostnader og Lobycos økonomiske bidrag. Denne mangelen betyr ikke at appen er uten verdi. Det betyr at rapportens anbefalinger bør handle om måling, governance og betinget investering, ikke om bastante lønnsomhetskonklusjoner.

Kildekritikken påvirker også hvordan Lobyco skal behandles i rapporten. Caseoppgaven åpner for spørsmålet om Coop bør selge eller afhende Lobyco. Det offentlige materialet som er gjennomgått, gir ikke et sterkt nok grunnlag til å anbefale salg som en evidensbasert konklusjon. Samtidig gir det heller ikke et sterkt nok grunnlag til å hevde at Lobyco nødvendigvis bør beholdes uendret. Den mest forsvarlige bruken av evidensen er å behandle Lobyco som et governance- og kapabilitetsspørsmål: Coop bør vurdere Lobycos mandat, kostnader, ekstern verdi, avhengighet til Coop App og bidrag til Coops butikkbaserte forretningsmodell før direksjonen konkluderer.

Kildekritikken skjerper også retail media-spørsmålet. Retail media kan være en mulig value capture-mekanisme fordi app, medlemsdata, leverandørkampanjer og butikkinnløsning kan kobles sammen. Men det er utrygt å skrive at retail media alene bærer appinvesteringen eller at supplier-funded campaigns dokumenterer nettoeffekt. Rapporten bør i stedet formulere retail media som et mulig bidrag som må testes med tydelige KPI-er, kostnadsdata og tillitsvurderinger. For en medlemseid aktør som Coop er kundetillit ikke bare et kommunikasjonsproblem, men en del av forretningsmodellens legitimitet.

Denne kildekritiske disiplinen er særlig viktig for Martin Hasgard Olesen. Som direktør for kommunikasjon, marketing og digital vil han naturlig møte argumenter om reach, engagement, kampanjer, personalisering og medlemskontakt. Hans bidrag i direksjonsdiskusjonen bør ikke være å avvise disse verdiene, men å oversette dem til et mer presist styringsspråk. Han bør kunne si hvilke appmålinger som viser digital aktivitet, hvilke som viser kundeverdi, og hvilke som faktisk kan knyttes til butikkøkonomi eller value capture. Dermed blir kildekritikken ikke bare en metodeøvelse, men en del av selve strategien.

\section{Teorivalg og anvendelse}

Oppgaven bruker DVC Framework og Digital Business Model-perspektivet som de eneste fulle teoriperspektivene. Valget er gjort fordi de to perspektivene utfyller hverandre i Coop-casen. DVC hjelper med å analysere hva slags digital verdi Coop App kan skape for kunder, medlemmer og relasjoner. Digital Business Model-perspektivet hjelper med å vurdere om denne verdien faktisk passer inn i Coops butikkbaserte forretningsmodell og kan omsettes til styrbar value creation og value capture.

Denne teoribruken skal være anvendt, ikke generell. Oppgaven skal derfor ikke bruke lange teorikapitler om digital transformasjon, plattformer eller konkurransestrategi. Den sentrale logikken er:

\begin{quote}
Teoribegrep $\rightarrow$ Coop-observasjon $\rightarrow$ analytisk betydning $\rightarrow$ strategisk implikasjon
\end{quote}

\subsection{DVC Framework}

DVC Framework brukes selektivt gjennom fem begreper: experiences, relationships, relevance, digital infrastructure og digital outputs. Experiences er relevant fordi Coop App kan påvirke selve handleopplevelsen gjennom Scan \& Betal, tilbud, kvitteringer, handleliste, spill og butikkkommunikasjon. Den analytiske betydningen er at appfunksjoner bør vurderes etter om de løser reelle kunde- eller butikkproblemer, ikke bare etter om de skaper aktivitet. Strategisk impliserer dette at Coop bør prioritere appfunksjoner som gjør butikkreisen enklere, mer relevant eller mer tillitsskapende.

Relationships er relevant fordi Coop ikke bare har kunder, men medlemmer og medeiere. Coop App kan styrke relasjonen gjennom personaliserte tilbud, medlemsfordeler, lokal kommunikasjon og gjentatte kontaktpunkter. Samtidig viser kildekritikken at relasjon ikke automatisk er det samme som lønnsomhet. Den analytiske betydningen er at appbasert lojalitet må skilles fra inkrementell verdi. Strategisk betyr dette at Martin bør forsvare medlemsrelasjonen som et viktig DVC-element, men samtidig kreve målinger som viser om relasjonen faktisk gir ønsket butikk- og forretningsmodellbidrag.

Relevance er særlig viktig fordi Coop har flere kjeder med ulike kundeløfter. En felles app kan gi skala og gjenkjennelighet, men den kan også bli for generell. 365discount-kunder kan ha andre digitale behov enn kunder i SuperBrugsen, Kvickly eller Brugsen. DVC-begrepet relevance gjør det mulig å analysere ``go deep''-spørsmålet uten å hoppe direkte til anbefalingen om separate apper. Strategisk peker dette mot kjedespesifikke value propositions innenfor én styrt plattform.

Digital infrastructure brukes til å forstå Lobyco. Lobyco kan være mer enn en leverandør eller et appteam; det kan være en digital infrastruktur for lojalitet, apputvikling, kampanjer og mulig retail media. Men DVC alene kan gjøre infrastrukturen for positivt ladet dersom den ikke kobles til økonomi. Derfor må Lobyco analyseres som mulig kapabilitet, men også som et kostnads-, governance- og prioriteringsspørsmål.

Digital outputs er det mest kildekritiske DVC-begrepet i denne casen. Brukere, monthly active users, household reach, spilldeltakelse, premier og engagement kan være viktige output-mål. Likevel er de mellomliggende mål. De viser at noe skjer digitalt, men ikke nødvendigvis at Coop skaper profitabilitet eller strategisk fordel. Strategisk betyr dette at rapporten senere bør anbefale et tydelig KPI-hierarki som skiller mellom adoption, engagement, kundeatferd, butikkøkonomi og value capture.

\subsection{Digital Business Model-perspektivet}

Digital Business Model-perspektivet brukes til å teste om DVC-verdiene faktisk passer Coops forretningsmodell. Det første relevante begrepet er value proposition. Coop App kan tilby sparing, bekvemmelighet, personalisering, medlemsfordeler, klima- og produktinformasjon, spill og lokal butikkrelevans. Dette er en bred value proposition. I en kjedeportefølje kan bredden være en fordel, men den kan også gjøre appens strategiske rolle uklar. Den analytiske oppgaven blir å spørre hvilken value proposition appen egentlig skal ha for hver kjede.

Value creation handler om hvem som får verdi og hvordan. Coop App kan skape verdi for kunder gjennom tilbud og bekvemmelighet, for butikker gjennom trafikk og bedre kampanjeaktivering, for leverandører gjennom kampanjeflater, og for Coop gjennom data, medlemskontakt og mulig retail media. Lobyco kan i tillegg skape verdi som kompetanse- og teknologikapabilitet. Men flere mulige verdiområder gjør også styringen vanskeligere. En app som forsøker å skape verdi for alle parter samtidig, trenger klare prioriteringer.

Value capture er det mest kritiske punktet. Den offentlige evidensen kan peke på mulige mekanismer som butikktrafikk, kurvstørrelse, handlefrekvens, supplier-funded campaigns, retail media, retention og kundetillit. Den dokumenterer likevel ikke netto lønnsomhet. Digital Business Model-perspektivet gjør derfor analysen mer disiplinert enn et rent digital value-perspektiv: det er ikke nok at appen skaper opplevelse, relasjon eller engagement; Coop må kunne fange verdi på en måte som passer økonomien, kjedene og butikkmodellen.

Key resources and capabilities brukes til å analysere Lobyco, medlemsbasen, appdata, butikknettverket, private labels og leverandørrelasjoner. Coop kan ha en interessant kombinasjon av fysisk skala og digital kundekontakt. Spørsmålet er om denne kombinasjonen styres som en Coop-first kapabilitet, eller om den blir en kostbar digital aktivitet med svak kobling til butikkøkonomien. Her blir hybrid governance relevant som senere stance: Lobyco bør verken idealiseres som automatisk konkurransefortrinn eller avskrives som distraksjon uten bedre økonomisk evidens.

Channels-begrepet brukes til å plassere appen riktig. Etter lukking av Coop.dk MAD og Coop.dk-webshop bør Coop App analyseres som en butikkstøttende kanal, ikke som en erstatning for en full online grocery-kanal. Appens strategiske verdi ligger derfor i koblingen mellom digital kontakt og fysisk butikkreise: tilbud som brukes i butikk, Scan \& Betal, lokale butikkvalg, medlemsfordeler, kvitteringer og kampanjeinnløsning.

Til slutt brukes governance og KPI-er som et eksplisitt business model-problem. Dersom Coop må velge mellom butikkforbedringer, prisinvesteringer, logistikk, kjedeposisjonering, apputvikling og Lobyco-ambisjoner, kan ikke digital suksess defineres bredt som ``mer engagement''. Martin bør derfor senere anbefale en KPI-logikk som starter med digitale output-mål, men ikke stopper der. Den bør koble appen til store traffic, incremental basket, margin, kjedespesifikk relevans, supplier-funded contribution, kundetillit og kostnad per funksjon.

Overgangen til den strategiske situasjonsanalysen blir dermed slik: DVC viser hvorfor Coop App kan være verdifull som kundeopplevelse, relasjonskanal, relevansmotor, digital infrastruktur og output-system. Digital Business Model-perspektivet tester om disse verdiene faktisk passer Coops butikkbaserte modell, økonomiske situasjon og governance-behov. Den videre analysen bør derfor undersøke appen og Lobyco som betingede strategiske aktiva: verdifulle dersom de styrker butikkøkonomi, kjedefit og medlemsrelasjon; problematiske dersom de først og fremst skaper aktivitet, kompleksitet og leverandørclaims.

\todo{Sett inn presise teori-referanser til DVC Framework og Digital Business Model-pensum når referanselisten bygges. Ikke introduser flere fulle teoriperspektiver uten at teorivalget åpnes på nytt.}

\section{Digital Business Model-analyse: Coop som store-first omnichannel-aktør}

\subsection{Plassering av Coop i DBM-matrisen}

Weill og Woerner organiserer digital business model-analysen rundt to akser: graden av \textit{customer knowledge} virksomheten besitter, og om forretningsdesignet opererer innenfor en kontrollert verdikjede eller som del av et bredere ecosystem. Plasseringen i denne matrisen er ikke en teknisk klassifisering; den er en strategisk vurdering av hvor verdiskaping og verdifangst faktisk skjer.

For Coop er svaret tydelig: selskapet bør forstås som en \textit{store-first omnichannel}-aktør. Kjerneøkonomien ligger i de fysiske butikkene, og 2025-tallene understreker dette. Med en nettoomsetning på om lag 32,5~mrd.~DKK gjennom Coop~Danmark~A/S og 900 butikker inkludert brugsforeningene, er den fysiske handelen ikke bare en salgskanal --- den er forretningsmodellen. Coop~App binder et lag av digital kapabilitet rundt denne kjernen: Scan~\&~Betal, personlige tilbud, bonus, digitale kvitteringer, handlelister og lokal butikkkommunikasjon er alle funksjoner som eksisterer for å gjøre det fysiske butikkbesøket mer relevant, mer sømløst og mer datadrevet. Appen er derfor ikke i første rekke en selvstendig digital salgskanal. Den kan forstås som en strategisk kapabilitet som støtter butikkmodellen ved å koble medlemsrelasjon, kjøpsdata og kundeopplevelse tettere til det fysiske handlestedet.

Det er verdt å skille mellom to mulige tilleggsdimensjoner. Lobyco --- Coops heleide datterselskap som drifter appinfrastrukturen og har begynt å selge plattformen til andre aktører --- kan tolkes som en mulig \textit{modular producer}-kapabilitet i Weill og Woerners terminologi: en leverandør av gjenbrukbar lojalitets- og betalingsinfrastruktur som kan plugges inn i andres systemer. Dette er en reell opsjon, men den bør holdes analytisk adskilt fra Coops hovedmodell. At Lobyco selger plattformen til eksterne aktører, endrer ikke det faktum at Coops konkurranseevne avhenger av fysisk handel, ikke av B2B-plattformvirksomhet. \textit{Ecosystem driver}-rollen --- der Coop ville bli kundenes primære destinasjon for planlegging av hverdagen rundt dagligvarer --- er en langsiktig opsjon som ikke kan underbygges av tilgjengelig evidens på dette tidspunktet. Å behandle Coop som ecosystem driver i dag ville innebære å overstige det kildematerialet faktisk støtter.

\subsection{Hvorfor appen er strategisk viktig --- og hva den koster}

Weill og Woerner hevder at alle virksomheter bør begynne med å kartlegge sin digitale trussel og mulighet presist, fordi en feildiagnosert trussel leder til feil strategisk respons. For Coop er det avgjørende at denne trusselen ikke snevres inn til spørsmålet om online grocery. Lukkingen av Coop.dk-webshop i januar 2025 viser at direkte netthandel i mat ikke er Coops vei fremover. Men det betyr ikke at det digitale presset er borte.

Den egentlige trusselen er av en annen karakter: \textit{kunderelasjonen rundt dagligvarehandelen digitaliseres}, uavhengig av om selve kjøpet skjer fysisk. Kunder bruker i stigende grad apper og digitale flater til å planlegge handleturer, sammenligne priser, finne tilbud og opprettholde lojalitet til bestemte kjeder. Dersom denne digitale kunderelasjonen eies av en konkurrent eller en tredjepartsplattform, kan Coop miste kunnskapen om og relevansen overfor sine egne kunder --- selv om kundene fortsetter å handle i Coops butikker. Lavprisaktører som Netto og Lidl har vist at kuponger, member prices og sparelogikk digitalt kan forsterke en prisbevisst posisjon. Konkurrenter som Salling Group har Salling-appen med funksjoner som i stor grad overlapper med Coop App. For Coop, som allerede opererer under margin- og resultatpress, indikerer dette at det digitale ikke er et nice-to-have: det er et forsvarsverktøy for kunderelasjonen.

Mulighetsbildet er likevel reelt. Coops styrker er konkrete: en appbase på om lag 1,8~millioner brukere ifølge Lobyco-materiale, et etablert kooperativt medlemskap, fire kjeder med ulik posisjonering, et landsdekkende butikknett, egne merkevarer og en kjøpsdatabase som kan gi analytisk innsikt. Dersom disse ressursene utnyttes systematisk, kan Coop~App bidra til å øke relevant handlefrekvens, styrke kjedepreferanse og drive mer presise tilbud enn konkurrentene. Det er nettopp dette Weill og Woerner kaller \textit{customer knowledge}: evnen til å kjenne kunden dypere enn konkurrentene og omsette den kunnskapen til relevante opplevelser og tilbud.

Objeksjonen er imidlertid like viktig: muligheter koster penger. Coop rapporterte et EBIT på minus 215~mill.~DKK for 2025. I en slik situasjon kan ikke appinvesteringer behandles som strategiske ønskelister. De må vurderes opp mot alternativkostnad, tilgjengelig kapital og finansiell realisme. Strategisk ambisjon uten budsjettkontroll er ikke strategi. Dette betyr ikke at appinvesteringer bør stoppes, men at de bør prioriteres strengt og knyttes til dokumentert effekt på kjerneøkonomien.

\subsection{Hva Coop kan bygge videre på: innhold, opplevelse og plattform}

Weill og Woerner peker på tre dimensjoner der virksomheter kan bygge digital konkurransefordel: \textit{content}, \textit{customer experience} og \textit{platform}. For Coop er alle tre tilgjengelige, men de krever prioritering og tydelig kjedespesifikk differensiering.

\textit{Content} handler om hvilken informasjon, hvilke produkter og hvilken relevans Coop tilbyr kundene digitalt. Coop har sterk content-kapabilitet: egne merkevarer, lokale tilbud, tilbudsaviser, oppskrifter, klima- og bærekraftsinnsikt og kjedespesifikk kommunikasjon. Denne kapitalen er relativt unik, særlig kombinasjonen av lokal butikkrelevans og et kooperativt verdigrunnlag. Utfordringen er at mye av dette innholdet i dag kommuniseres flatt, uten tilpasning til den individuelle kundens handlehistorikk eller preferanser.

\textit{Customer experience} er etter alt å dømme Coop~Apps sterkeste strategiske bidrag i dag. Scan~\&~Betal fjerner kø og forenkler betaling. Digitale kvitteringer erstatter papirkvitteringer og gir grunnlag for kjøpsanalyse. Handlelisten integrerer planlegging med butikkbesøket. Bonus og personlige tilbud gjør hvert handlebesøk mer relevant. Disse funksjonene gjør ikke dagligvarehandel digital i seg selv --- de gjør det \textit{lettere og mer relevant å handle i Coops butikker}. Det er den rette strategiske logikken for en store-first omnichannel-aktør.

\textit{Platform} dreier seg om evnen til å samle appopplevelse, data og verdistrøm på én felles infrastruktur som kan skaleres. Her er Lobycos rolle relevant: én felles digital plattform på tvers av Coops kjeder kan gi datakvalitet, kostnadseffektivitet og gjenbruk av funksjonalitet. Men plattformlogikken må kombineres med kjedespesifikk differensiering. 365discount-kunder trenger sparelogikk, prissammenligning og kuponger; Brugsen-kunden søker lokal relevans og kjennskap til nærbutikken; SuperBrugsen-kunden verdsetter inspirasjon og kvalitetsrelevans; Kvickly-kunden planlegger bredere handlebesøk og trenger oversikt og effektivitet. En ren one-size-fits-all-app risikerer å levere middelmådige opplevelser til alle disse kundegruppene. En analytisk presis formulering for Coops plattformstrategi er: \textit{én felles plattforminfrastruktur, kjedespesifikke appmoduler og felles KPI-governance}.

\subsection{Strategisk spenning: store-first omnichannel versus Lobyco som modular producer}

Den sentrale strategiske spenningen i Coop~App-casen kan formuleres slik: skal Lobyco primært være en kapabilitet \textit{for} Coop, eller kan den over tid bli et selvstendig plattformselskap med egen vekstlogikk?

Argumentet for modular producer-rollen er teknologisk og finansielt interessant. Lobyco har investert tungt i lojalitets- og appinfrastruktur, og dersom denne infrastrukturen kan selges til andre dagligvareaktører i Europa, kan det gi skalerte inntekter uten å begrense Coops egen appbruk. Retail media-inntekter fra leverandørfinansierte kampanjer er et annet eksempel på at Coop~App-plattformen kan generere verdier utover direkte salg.

Risikoen er imidlertid reell og bør ikke undervurderes. Dersom Lobycos eksterne kunder og vekstambisjoner begynner å definere produktveikart, prioriteringer og investeringsbehov, kan fokus flyttes bort fra Coops kjerneforretning. Det er en organisatorisk logikk i plattformselskaper der vekst og skalering til slutt prioriteres over vertsvirksomhetens spesifikke behov. For Coop, som allerede er under finansielt press, er risikoen at Lobyco absorberer ledelsesoppmerksomhet og kapital uten at Coop-verdien er tilstrekkelig dokumentert.

Den analytiske konklusjonen er at Lobyco bør forstås og styres som en \textit{Coop-first modular producer}: plattformkapabiliteten er reell og verdifull, men skalering til eksterne aktører bør kreve at Lobycos bidrag til Coops egen kjerneforretning er demonstrert og dokumentert. KPI-governance er her avgjørende. Output-metrikker som månedlige aktive brukere, kampanjeklikk og appinstallasjoner er ikke tilstrekkelige alene. De forteller om aktivitet, ikke om verdi. Outcome-KPI-er som kobler appbruk til butikktrafikk, gjennomsnittlig handlekurv, kjedemargin, kundelojalitet over tid og leverandørbidrag er nødvendige for å vurdere om Coop~App faktisk skaper den verdien den bør skape for en store-first omnichannel-aktør.

Det er også analytisk relevant å trekke inn Coops historiske erfaringer med teknologisk gjeld. Bjørn-Andersen og Hedman dokumenterte i 2016 at Coop over tid hadde akkumulert IT-kompleksitet og styringsutfordringer som begrenset handlefriheten i strategiske valg. Poenget er ikke at denne situasjonen nødvendigvis er uendret i dag, men at digitale valg --- herunder apparkitektur, plattformavhengigheter og leverandørrelasjoner --- kan skape fremtidig kompleksitet og gjeld dersom de ikke styres bevisst. Digital strategi handler ikke bare om å legge til nye funksjoner, men om å styre kompleksitet og avhengigheter slik at fremtidige valg forblir åpne.

\subsection{Delkonklusjon: Coop~App som business model-komponent}

Den analytiske konklusjonen fra DBM-perspektivet er at Coop bør prioritere \textit{store-first omnichannel} som sin strategiske hovedretning, med Lobyco som en kontrollert modular producer-opsjon underlagt klar Coop-first-styring. Coop~App bør ikke vurderes som en isolert applikasjon med egne suksesskriterier, men som en \textit{business model-komponent} der verdien måles i bidrag til Coops kjerneøkonomi: butikktrafikk, margin, kundelojalitet og kjederelevans.

Det er i Weill og Woerners begrep \textit{customer knowledge} den strategiske verdien av Coop~App kan bli størst. Coop har et unikt utgangspunkt: et etablert kooperativt medeierskapsforhold, et appbasert kjøpshistorikk-grunnlag og et landsdekkende butikknett som gir konkret kontekst for databruk. Dersom Coop klarer å omsette denne kunnskapen til mer presise, relevante og differensierte opplevelser enn konkurrentene, kan appen styrke Coops posisjon på lang sikt. Men det er et \textit{dersom}: appbruk er ikke det samme som kundekunnskap, og kundekunnskap er ikke det samme som verdiskapende beslutninger. Veien fra data til verdi krever analytisk kapasitet, organisatorisk evne og en governance-struktur som kobler digitale investeringer til forretningsresultater.

Videre skalering av digitale satsinger --- enten det gjelder nye appfunksjoner, Lobycos vekst eller retail media-ambisjoner --- bør derfor kreve dokumentert Coop-verdi som forutsetning. Digitaliseringen skal støtte og forsterke Coops kjerneforretning i de fysiske butikkene, ikke skape et parallelt digitalt spor løsrevet fra der kundene og verdiene faktisk befinner seg. Dette leder naturlig videre til den strategiske situasjonsanalysen og spørsmål om KPI-governance, kjedespesifikke moduler og betingede investeringskriterier.

\section{Strategisk situasjonsanalyse}

Coops strategiske situasjon i 2026 er preget av en spenning mellom en synlig digital kundekanal og en økonomisk presset, fysisk butikkbasert kjerne. Coop App og Lobyco kan representere viktige digitale kapabiliteter, men den strategiske verdien kan ikke vurderes isolert fra Coops butikknettverk, kjedeportefølje, medlemsmodell og økonomiske handlingsrom. Dette gjør appspørsmålet til mer enn et teknologivalg. Det er et spørsmål om hvilke digitale aktiviteter som faktisk støtter Coops forretningsmodell.

Det første analytiske poenget er at app- og Lobyco-beslutningen må forstås som et investeringsallokeringsproblem. Coop rapporterer positiv EBITDA før eiendomsgevinster, men også negativ EBIT. Dette indikerer at konsernet ikke har ubegrenset rom for brede digitale satsinger som ikke kan kobles til tydelig verdi. Den relevante vurderingen er derfor ikke om Coop App har mange brukere, men om videre apputvikling gir bedre bruk av knappe ressurser enn alternative investeringer i pris, butikker, logistikk, kjedeutvikling eller driftsforbedring.

DVC bidrar her ved å skille mellom digitale output-mål og dypere verdi. Appbruk, kampanjeaktivitet og engagement kan vise at Coop har digital kontakt med kundene. Men disse målene er bare mellomliggende indikatorer. Digital Business Model-perspektivet gjør spørsmålet strengere: Coop må vurdere om appaktivitet kan omsettes til value capture, for eksempel gjennom relevant butikktrafikk, inkrementell kurv, bedre margin, lavere friksjon, supplier-funded contribution eller sterkere medlemsretensjon. Uten en slik kobling blir appen en kostnads- og kompleksitetsdriver snarere enn et dokumentert strategisk aktivum.

Det andre analytiske poenget er at Coops digitale strategi må vurderes ut fra fysisk butikk som hovedkanal. Lukkingene av Coop.dk MAD og Coop.dk-webshop peker ikke på at digitalisering er irrelevant, men de svekker en ren e-commerce-tolkning av casen. Coop App bør derfor analyseres som et digitalt lag rundt butikkhandelen. Appen kan støtte planlegging, tilbud, betaling, kvitteringer, lokale butikkvalg, medlemsfordeler og kampanjeinnløsning. Den bør ikke primært forstås som en vei tilbake til en bred online grocery-modell.

Dette gir en klar strategisk implikasjon: Appens suksess må måles gjennom dens støtte til butikkøkonomien. Dersom Scan \& Betal reduserer friksjon i butikk, dersom personlige tilbud øker relevant innløsning, eller dersom kampanjer trekker kunder til butikk på en lønnsom måte, kan appen ha forretningsmodellverdi. Dersom funksjonene først og fremst skaper klikk, spilldeltakelse eller synlig digital aktivitet, er verdien langt mer usikker.

Det tredje analytiske poenget er at appreach ikke kan forveksles med strategisk styrke. Lobyco oppgir betydelig reach for Coop App, inkludert mer enn 1,8 millioner brukere og omtrent 25 prosent household reach. Dette er strategisk relevant fordi det indikerer at Coop har en stor digital kontaktflate. En app med høy reach kan være vanskelig å erstatte, særlig for en medlemsorganisasjon som trenger direkte kundekontakt. Men reach er ikke det samme som lønnsomhet. Den dokumenterer ikke kausal handlefrekvens, marginbidrag eller varig konkurransefortrinn.

Dette er en sentral kildekritisk og strategisk skillelinje. En optimistisk appfortelling vil kunne si at høy adoption viser at Coop bør investere mer. En pessimistisk fortelling vil kunne si at økonomisk press viser at Coop bør kutte. Begge tolkninger hopper over hovedspørsmålet: hvilken type verdi skaper appen, for hvem, til hvilken kostnad, og med hvilken dokumentert effekt? Rapportens videre anbefalinger må derfor bygge på et KPI-hierarki, ikke på ett enkelt suksessmål.

Det fjerde analytiske poenget gjelder kjedeforskjeller. Coop består av flere kjeder med ulike kundeløfter. 365discount, Brugsen, SuperBrugsen og Kvickly kan ikke uten videre antas å ha samme digitale behov. En discountkunde kan prioritere enkelhet, pris og raske tilbud, mens kunder i andre kjeder kan respondere mer på inspirasjon, ferskvarer, lokal tilhørighet, medlemsfordeler eller bredere butikkopplevelse. DVC-begrepet relevance blir derfor sentralt: digital verdi oppstår ikke bare fordi en funksjon finnes, men fordi den passer kundens situasjon.

Konkurrentevidensen styrker denne tolkningen, men må brukes forsiktig. Salling, REMA, Lidl, Dagrofa/MENY og Nemlig viser ikke at apper i seg selv skaper lønnsomhet. Den tryggere inferensen er at digitale løsninger ser ut til å være sterkest når de passer med aktørens forretningsmodell og kundeløfte. Lidl Plus kan forstås som tett koblet til rabatt- og kuponglogikk. REMA fremstår mer fokusert på enkel butikkutility. Salling har digitale løsninger som kan kobles til butikkflyt og skala. Dagrofa/MENY peker mot mat- og lokal identitet. Nemlig viser at online grocery krever en egen logistikk- og økonomimodell. For Coop betyr dette at ``go deep''-spørsmålet bør forstås som et spørsmål om kjedespesifikk relevans, ikke nødvendigvis som et krav om separate apper.

Det femte analytiske poenget gjelder Lobyco. Lobyco bør ikke behandles som et enkelt enten/eller-spørsmål. På den ene siden kan Lobyco representere en digital kapabilitet: apputvikling, lojalitetsinfrastruktur, data, kampanjelogikk og mulig retail media. På den andre siden mangler det offentlige kildegrunnlaget dokumentasjon på Lobycos kostnader, eksterne inntekter, avhengighet til Coop App og eventuell salgsverdi. Dette gjør det uansvarlig å konkludere bastant med at Lobyco bør selges, beholdes uendret eller styrkes kraftig.

Digital Business Model-perspektivet gjør Lobyco til et governance-spørsmål. Hvis Lobyco bidrar til Coops value proposition og value capture i butikkene, kan det være en strategisk kapabilitet. Hvis Lobyco utvikler seg i retning av en plattform- eller OEM-logikk som trekker ledelsesoppmerksomhet og investering bort fra Coops kjerne, kan det bli en distraksjon. Den relevante direksjonsdiskusjonen bør derfor ikke begynne med ``selg eller behold'', men med spørsmål om mandat, målinger, kostnader, kontroll og Coop-first-prioritering.

Det sjette analytiske poenget gjelder retail media og gamification. Kildene indikerer at appen kan brukes til leverandørfinansierte kampanjer, spill, premier og butikkinnløsning. Dette kan være en mulig value capture-mekanisme. Det kan også støtte DVC-dimensjonen experiences ved å skape aktivitet rundt butikk og medlemskap. Men kildene er i hovedsak leverandør- eller marketingkilder. De dokumenterer ikke nettoeffekt etter kampanjekostnader, premieutgifter, marginpåvirkning og mulig slitasje på kundetillit.

Derfor bør gamification og retail media behandles som testbare mekanismer, ikke som beviste strategiske svar. Hvis spill og kampanjer skaper butikktrafikk, relevant innløsning og supplier-funded contribution etter kostnader, kan de inngå i appens forretningsmodell. Hvis de først og fremst skaper lavverdi-engagement, kompleksitet eller en følelse av at medlemsdata brukes for aggressivt, kan de svekke tilliten. For en medlemseid aktør er dette viktig: kundetillit er ikke bare et kommunikasjonsmål, men en del av verdiforslaget.

Samlet viser analysen at Coop App og Lobyco hverken bør avskrives eller idealiseres. Coop har en stor fysisk og medlemsbasert kjerne, en betydelig digital kontaktflate og en mulig teknologisk kapabilitet. Samtidig er økonomien presset, kjedene forskjellige, kildegrunnlaget ujevnt og effekten av appbruk usikkert dokumentert. Det strategisk mest forsvarlige er derfor å styre appen og Lobyco som betingede aktiva: verdifulle dersom de kan knyttes til butikkverdi, kjederelevans, medlemsrelasjon og value capture; problematiske dersom de først og fremst produserer aktivitet, kompleksitet og leverandørclaims.

\section{Strategisk stance}

På bakgrunn av analysen bør Martin Hasgard Olesen argumentere for hybrid governance. Dette betyr at han ikke bør gå inn i direksjonsdiskusjonen med en enkel posisjon om å styrke appen, kutte appen eller selge Lobyco. Han bør i stedet argumentere for at Coop beholder Coop App og Lobyco som mulige strategiske kapabiliteter, men gjør videre investering betinget av dokumentert kunderelevans, butikkverdi, kjedespesifikk fit, tillit og klare KPI-er.

Denne posisjonen svarer på de viktigste alternativene i casen. En ren reduksjonsstrategi ville ta økonomisk press på alvor, men risikerer å svekke en digital kontaktflate og medlemsrelasjon som Coop allerede har bygget opp. En ren styrkingsstrategi ville ta appens reach på alvor, men risikerer å forveksle adoption og engagement med forretningsmessig effekt. Et salg eller en afhændelse av Lobyco kan være relevant dersom økonomisk og strategisk dokumentasjon viser at verdien er høyere utenfor Coop enn innenfor Coop, men dagens offentlige kildegrunnlag er ikke sterkt nok til en slik konklusjon. En full ``go deep''-strategi med separate eller sterkt divergerende apper kan gi bedre kjedefit, men kan også skape kostnader, kompleksitet og svekket governance.

Hybrid governance er derfor ikke et kompromiss uten retning. Det er en styringsposisjon som skjerper hva Coop må bevise før videre investering skaleres. Posisjonen kan formuleres slik: Coop bør beholde én styrt app- og lojalitetsplattform, men utvikle kjedespesifikke value propositions, målinger og moduler der det gir dokumentert verdi. Lobyco bør behandles som en Coop-first digital kapabilitet, ikke som et selvstendig prestisjeprosjekt. Retail media og gamification bør utvikles bare der de støtter kunderelevans, butikkverdi og tillit.

For Martin er denne stance særlig passende. Hans ansvarsområde ligger nær kommunikasjon, marketing, digital kundekontakt, lojalitet og apputvikling. Samtidig kan han ikke alene eie butikkøkonomi, margin eller driftsprioriteringer. Derfor bør hans rolle være å oversette digitale muligheter til direksjonens økonomiske og strategiske språk. Han bør vise hvordan appen kan bidra til butikkbasert forretningsmodell, men også insistere på at digitale claims må kunne testes.

\section{Strategiske anbefalinger}

De følgende anbefalingene følger av analysen. De er skrevet til Martin Hasgard Olesen og skal gjøre ham i stand til å delta i direksjonsdiskusjonen med en kildekritisk, men handlingsorientert posisjon.

\subsection{Redefiner Coop App-suksess rundt forretningsmodell-fit og butikkverdi}

Martin bør først anbefale at Coop redefinerer hva det betyr at Coop App er vellykket. Appen bør ikke primært vurderes etter brukertall, kampanjeaktivitet, spilldeltakelse, reach eller engagement. Disse målene er relevante, men de er ikke sluttmål. De bør plasseres nederst i et KPI-hierarki som leder videre til kundeadferd, butikkøkonomi og value capture.

Et mer relevant KPI-hierarki bør skille mellom fem nivåer. Første nivå er adoption og reach: hvor mange bruker appen, og hvor ofte? Andre nivå er engagement: hvilke funksjoner, tilbud og kampanjer brukes? Tredje nivå er kundeatferd: fører bruken til butikkbesøk, tilbudsinnløsning, retensjon eller mer relevant handlemønster? Fjerde nivå er butikkøkonomi: påvirkes kurv, margin, kategori, trafikk eller friksjon i butikk? Femte nivå er value capture: skaper appen dokumentert bidrag gjennom lønnsom kundeadferd, supplier-funded contribution, retail media, kostnadsreduksjon eller sterkere medlemsrelasjon?

Lobyco bør underlegges samme logikk. Det bør ikke vurderes bare etter teknologisk modenhet, appreach eller eksterne plattformmuligheter. Lobyco bør behandles som en Coop-first kapabilitet der mandat, kostnader og investeringer må begrunnes gjennom bidrag til Coops butikkbaserte forretningsmodell. Dersom Lobyco kan dokumentere at det styrker medlemsrelasjon, kjedefit, kampanjekvalitet, retail media eller butikkverdi, kan det forsvares som strategisk kapabilitet. Dersom dette ikke kan dokumenteres, bør direksjonen vurdere reduksjon, omprioritering eller alternative governance-modeller.

Retail media bør også inn i denne anbefalingen, men med tydelige guardrails. Det bør behandles som en mulig value capture-mekanisme, ikke som bevis på at appen er økonomisk vellykket. Martin bør anbefale at retail media-mål inkluderer netto bidrag etter kostnader, kunderelevans, tillit og påvirkning på medlemsopplevelse. Hvis leverandørfinansierte kampanjer gir kortsiktig inntekt, men svekker opplevd relevans eller tillit, kan de skade forretningsmodellen på lengre sikt.

\subsection{Utvikle kjedespesifikke app-value propositions innenfor én styrt Coop-plattform}

Den andre anbefalingen er at Coop bør utvikle kjedespesifikke app-value propositions, men innenfor én styrt Coop-plattform. Dette er rapportens svar på casealternativet om å ``go deep''. Coop bør ikke uten videre lage separate apper for hver kjede, fordi det kan øke kostnader, fragmentere data, skape teknisk kompleksitet og svekke felles governance. Samtidig bør Coop heller ikke anta at én bred appopplevelse er like relevant for alle kjeder.

For 365discount bør appverdien sannsynligvis vektlegge pris, enkelhet, tydelige tilbud, rask butikkreise og lav kompleksitet. For Brugsen kan lokal relevans, medlemsrelasjon og nærhet til butikk være mer sentralt. For SuperBrugsen og Kvickly kan inspirasjon, ferskvarer, bredde, personlig relevans og handleplanlegging være viktigere. Dette er foreløpige strategiske hypoteser, ikke dokumenterte kundeinnsikter. De bør testes med kjededata før skalering.

Anbefalingen er derfor modulær: én plattform, men differensierte flater, budskap, tilbudslogikk, kampanjer og KPI-er. DVC støtter dette gjennom relevance-begrepet: digital verdi må passe kundens situasjon. Digital Business Model-perspektivet støtter det gjennom value proposition og channels: appen bør ikke bare være en teknisk kanal, men en kanal som uttrykker hver kjedes kundeløfte og forretningsmodell.

Denne anbefalingen bør også redusere risikoen for at Coop overinvesterer i appfunksjoner som bare gir verdi for noen kjeder. Dersom en funksjon er viktig for SuperBrugsen, men lite relevant for 365discount, bør dette synliggjøres i governance og KPI-er. Coop bør derfor ikke spørre om en funksjon er ``god for appen'' generelt, men om den er god for en bestemt kjede, kundemission og butikkøkonomi.

\subsection{Gjennomgå og redesign gamification- og engagement-funksjoner}

Den tredje anbefalingen er at Coop bør gjennomgå og redesigne gamification- og engagement-funksjoner hvis de ikke kan kobles til kunde- eller butikkverdi. Dette er ikke en påstand om at gamification ikke virker. Det er en kildekritisk styringsanbefaling: engagement må vurderes som middel, ikke som mål.

Playable- og lignende kilder kan indikere at spill, premier og kampanjer skaper aktivitet og butikkinnløsning. Dette er relevant, særlig dersom kampanjene kan finansieres av leverandører og føre til butikktrafikk. Men kildene er leverandørkilder og dokumenterer ikke nødvendigvis nettoeffekt. Coop trenger derfor en feature review-modell som klassifiserer engagement-funksjoner i tre grupper.

Første gruppe er funksjoner som bør beholdes fordi de forbedrer butikkreisen, tilbudsrelevans, medlemsverdi eller dokumentert kampanjebidrag. Andre gruppe er funksjoner som bør redesignes fordi de skaper aktivitet, men ikke nok tydelig verdi. Tredje gruppe er funksjoner som bør fjernes eller nedprioriteres fordi de øker kompleksitet, svekker tillit eller ikke kan kobles til butikkøkonomi.

DVC hjelper med å analysere dette gjennom experiences og digital outputs. En god digital opplevelse kan skape verdi, men output-mål må ikke forveksles med outcome. Digital Business Model-perspektivet skjerper vurderingen ved å spørre om funksjonen bidrar til value creation og value capture. Martin bør derfor ikke avvise gamification, men kreve at den beviser sin rolle i en butikk- og medlemsorientert forretningsmodell.

\subsection{Bruk Coop App til å styrke medlemsrelasjonen, men skill lojalitet fra lønnsomhet}

Den fjerde anbefalingen er at Coop fortsatt bør bruke appen til å styrke medlemsrelasjonen, men at Martin må skille tydelig mellom lojalitet og lønnsomhet. Coop har en særskilt medlems- og medeieridentitet. Appen kan gjøre denne identiteten mer konkret gjennom personlige tilbud, digitale kvitteringer, lokal kommunikasjon, medlemsfordeler og oversikt over bonus eller kjøp. Dette kan være strategisk verdifullt også når verdien ikke umiddelbart lar seg måle i kortsiktig margin.

Samtidig er det farlig å gjøre medlemslojalitet til et lønnsomhetsbevis. Appbrukere kan allerede være mer lojale kunder. Hvis de handler oftere, kan det skyldes at de var mer Coop-orienterte før de begynte å bruke appen. Uten kontrollgrupper, før-/etter-data og margininformasjon kan ikke rapporten påstå at appen skaper profitabel lojalitet.

Martin bør derfor anbefale at Coop skiller mellom tre typer medlemsverdi. Den første er relasjonsverdi: appen gjør Coop mer relevant, tilgjengelig og personlig for medlemmet. Den andre er atferdsverdi: appen kan påvirke besøk, innløsning, handleplanlegging eller retention. Den tredje er økonomisk verdi: atferden gir faktisk bidrag etter kostnader og margin. Alle tre er viktige, men de bør ikke blandes.

Denne anbefalingen passer godt med DVC-begrepet relationships. Coop App kan være en relasjonskanal, ikke bare en rabattmaskin. Digital Business Model-perspektivet krever likevel at relasjonsverdien knyttes til value capture dersom den skal brukes som investeringsargument. For Martin er dette et viktig direksjonspoeng: Coop bør beskytte medlemsrelasjonen, men må samtidig måle hva den digitale relasjonen faktisk gjør for forretningsmodellen.

\subsection{Pilotér kjedenivå-initiativer før større appinvesteringer skaleres}

Den femte anbefalingen er at Coop bør pilotere kjedenivå-initiativer før større appinvesteringer skaleres på tvers av konsernet. Denne anbefalingen må skilles fra anbefaling 2. Anbefaling 2 handler om strategisk design: appen bør ha kjedespesifikke value propositions innenfor én plattform. Anbefaling 5 handler om implementering og risikoreduksjon: Coop bør teste disse value propositions før de gjøres til store, permanente investeringer.

Pilotene bør defineres med klare hypoteser. For eksempel kan Coop teste om en enklere discount-orientert appmodul øker relevant tilbudsinnløsning i 365discount uten å øke kompleksiteten. Coop kan teste om lokal butikkkommunikasjon gir sterkere medlemsrelasjon i Brugsen. Coop kan teste om inspirasjons- og ferskvarelogikk skaper relevant handleplanlegging i SuperBrugsen eller Kvickly. Poenget er ikke å bevise alt på forhånd, men å hindre at Coop skalerer brede appinvesteringer uten læring.

Hver pilot bør ha forhåndsdefinerte målinger. Minimum bør Coop måle adoption, aktiv bruk, tilbudsinnløsning, butikkbesøk, kurv og margin der data finnes, kundetillit, kostnad per funksjon og eventuelt supplier-funded contribution. Dersom en pilot bare skaper engagement, men ikke viser kunderelevans eller forretningsmodellbidrag, bør den ikke skaleres uten redesign.

Denne anbefalingen støtter hybrid governance. Den gjør det mulig å bevare appen og Lobyco som kapabiliteter uten å gi dem blankofullmakt. Den gir også Martin en konstruktiv rolle i direksjonen: han kan argumentere for digital utvikling, men på en måte som respekterer Coops økonomiske press og butikkbaserte kjerne.

\section{Begrensninger og kritisk refleksjon}

Rapportens viktigste begrensning er mangelen på interne Coop-data. Offentlige kilder kan dokumentere butikkstruktur, regnskapstall, appfunksjoner, medlemsmodell og leverandørenes egne claims. De kan derimot ikke dokumentere appens inkrementelle effekt på margin, kurv, butikkbesøk, kampanjelønnsomhet, kundetillit eller Lobycos økonomiske bidrag. Derfor er anbefalingene formulert som governance-, måle- og testanbefalinger, ikke som bastante konklusjoner om appens lønnsomhet.

En annen begrensning er kildegrunnlagets kommersielle bias. Lobyco, Playable og Shortcut er relevante fordi de kjenner appen, kampanjelogikken og utviklingshistorikken. Samtidig har de interesse av å fremstille appen og egne løsninger som vellykkede. Rapporten bruker derfor disse kildene som self-presentation, intended value og claims, ikke som nøytral effektmåling.

En tredje begrensning er selection bias i appbrukerclaims. Dersom appbrukere handler oftere, kan det skyldes at appen påvirker dem. Men det kan også skyldes at de mest lojale Coop-kundene er mest tilbøyelige til å bruke appen. Uten kontrollgrupper og før-/etter-data bør slike claims ikke brukes som kausal dokumentasjon.

Konkurrentevidensen har også begrensninger. Den viser mønstre og kontekst, men den beviser ikke at apper skaper konkurrentenes økonomiske resultater. Salling, REMA, Lidl, Dagrofa/MENY og Nemlig brukes derfor til å vise digital fit med forretningsmodell, ikke til å hevde at Coop bør kopiere bestemte funksjoner.

Til slutt har rapporten en formell språkbegrensning. Prosjektet og dette utkastet er skrevet på norsk etter brukerens valg. Dersom oppgaven formelt krever dansk eller engelsk ved innlevering, må språkvalget avklares før endelig hand-in.

\section{Konklusjon}

Martin Hasgard Olesen bør gå inn i direksjonsdiskusjonen med en hybrid governance-posisjon. Coop bør ikke behandle Coop App som et isolert digitalt prestisjeprosjekt, men heller ikke redusere eller selge digitale kapabiliteter uten en tydeligere vurdering av deres forretningsmodellbidrag. Appen og Lobyco bør forstås som betingede strategiske aktiva: verdifulle dersom de styrker butikkøkonomi, kjedespesifikk relevans, medlemsrelasjon og value capture; problematiske dersom de hovedsakelig skaper aktivitet, kompleksitet og leverandørclaims.

Analysen viser at Coop fortsatt er en fysisk butikkbasert dagligvareaktør, at appen bør vurderes som en butikkstøttende kanal, og at Lobyco må behandles som et governance- og kapabilitetsspørsmål. DVC Framework viser hvordan appen kan skape digital kundeverdi gjennom opplevelser, relasjoner, relevans, infrastruktur og output. Digital Business Model-perspektivet viser at denne verdien først blir strategisk sterk når den passer Coops forretningsmodell og kan kobles til value creation og value capture.

De fem anbefalingene følger av denne logikken. Coop bør redefinere app-suksess rundt forretningsmodell-fit og butikkverdi, utvikle kjedespesifikke app-value propositions innenfor én styrt plattform, gjennomgå gamification og engagement-funksjoner kildekritisk, bruke appen til å styrke medlemsrelasjonen uten å forveksle lojalitet med lønnsomhet, og pilotere kjedenivå-initiativer før store investeringer skaleres. Slik kan Martin forsvare digital utvikling, men samtidig stille krav til måling, prioritering og strategisk disiplin.

\section{Referanser}

\todo{Bygg referanseliste i systematisk stil, fortrinnsvis APA 7 dersom ikke annet bestemmes. Ikke finn på sidetall, forfattere eller titler før referansene er verifisert.}

Foreløpige referansebehov:

\begin{itemize}
  \item Coop Danmark A/S årsrapport 2025.
  \item Coop pressemelding om 2025-resultat.
  \item Coop kundeservice/side om lukking av Coop.dk-webshop.
  \item Coop App- og medlemssider.
  \item Lobyco Coop case og Lobyco insight om multi-banner loyalty.
  \item Playable Coop case.
  \item Shortcut Coop Denmark App case.
  \item Salling Group key figures og relevante app-/Scan\&Go-sider.
  \item REMA, Lidl, Dagrofa/MENY og Nemlig-kilder brukt i konkurrentkontekst.
  \item DVC Framework-pensum.
  \item Digital Business Model-pensum.
\end{itemize}

\section{AI-bruk}

\todo{Tilpass endelig tekst til CBS-/kursregler før innlevering.}

Generativ AI er brukt til å strukturere prosjektarbeidet, etablere agentroller, sammenstille kildekritiske arbeidsnotater og skrive arbeidsutkast. Brukeren har låst sentrale valg om rådgivningsaktør, teoriperspektiver, strategisk stance, språk og anbefalingsretning. AI-bruken i dette utkastet er synliggjort gjennom agentmodellen, faktaverifikasjonstrinnet og TODO-markører der kilder eller presise referanser må kontrolleres videre.

\end{document}
